\section{Niedervolt-Versorgungsspannungen kontrollieren}

Um sicherzustellen, dass alle Module in den NIM-Reihen mit den korrekten Spannungen versorgt wurden, wurden mit einer Probe, die an CH1 des Oszilloskops angeschlossen war, die anliegenden Versorgungsspannungen kontrolliert \cite[3]{525va}. Das Oszilloskop wurde statt einem Digital-Multimeter verwendet, um zeitliche Schwankungen erkennen zu können. Hierbei befanden sich die Module wie in Abbildung \ref{fig:aufbau_crates} dargestellt \todo{ref korrekt mit link}. Die für die einzelnen möglichen Versorgungsspannungen ($\SI{-6}{\V}, \SI{6}{\V}, \SI{-12}{\V}, \SI{12}{\V}, \SI{-24}{\V}, \SI{24}{\V}$ sowie die Erde, jeweils für die beiden Reihen getrennt) beobachteten Oszilloskopaufnahmen wurden als Spannung abhängig von der Zeit aufgetragen und sind in Abbildungen \ref{fig:crates_upper}, \ref{fig_crates_broken} und \ref{fig:crates_lower} dargestellt. Es zeigte sich bei der oberen Reihe, dass die Spannungen $\SI{-6}{\V}$ und $\SI{6}{\V}$ nicht der erwarteten Gleichspannung bei dem beschrifteten Spannungswert entsprachen, sondern eine höhere Amplitude und eine sinusförmige zeitliche Schwankung aufwiesen. Sie sind in Abbildung \ref{crates_broken} getrennt von den restlichen Versorgungsspannungen der Reihe aufgetragen, welche alle wie erwartet aussahen. Bei diesen beiden Spannungen war vermutlich der Gleichrichter und auch die Amplitudeneinstellung defekt, was für die Messungen im folgenden aber nicht relevant war, da sie nie als Versorgungsspannung verwendet wurden \todo{citation dafür mit manuals}. Die beobachteten Spannungen für die untere Reihe entsprachen in ihrem Verlauf alle den Erwartungen. Die notwendigen Versorgungsspannungen der NIMs waren damit alle funktionsfähig.

\jafps{crates_upper}{\todo{Caption fehlt}}{0.6}
\jafps{crates_broken}{\todo{Caption fehlt}}{0.6}
\jafps{crates_lower}{\todo{Caption fehlt}}{0.6}
